\lecture{15}{Функции Ляпунова в оптимизации}

\sourcevideo
    {https://www.youtube.com/playlist?list=PLOzRYVm0a65fhKDS8cYIC7YCuVGJ4FOJx}
    {Week 4: Lecture 15: Connections to Optimization Problems}

Функция Ляпунова --- скалярная характеристика состояния, которая
принимает минимальное значение на исследуемом равновесии и не
возрастает вдоль траекторий. Скорость её убывания позволяет доказать
устойчивость, асимптотическую или экспоненциальную сходимость. В
оптимизации естественными кандидатами служат расстояние до множества
решений, ошибка по значению функции и норма градиента.

\section{Классы функций и определённость}

\begin{definition}[Функция класса \(\mathcal K\)]
Непрерывная функция
\[
    \alpha\colon[0,a)\to[0,\infty)
\]
принадлежит классу \(\mathcal K\), если
\[
    \alpha(0)=0
\]
и \(\alpha\) строго возрастает. Если область определения равна
\([0,\infty)\) и
\[
    \alpha(r)\to\infty
    \qquad
    \text{при }r\to\infty,
\]
то \(\alpha\in\mathcal K_\infty\).
\end{definition}

Функции \(r\), \(r^2\) и \(r^p\) при \(p>0\) принадлежат
\(\mathcal K_\infty\). Функция
\[
    \alpha(r)=1-e^{-r}
\]
принадлежит \(\mathcal K\), но не \(\mathcal K_\infty\), поскольку
ограничена сверху единицей.

Пусть \(V\colon\RR^n\to\RR\) непрерывна. Она называется положительно
полуопределённой относительно начала координат, если
\[
    V(0)=0,
    \qquad
    V(x)\ge0,
\]
и положительно определённой, если дополнительно
\[
    V(x)>0
    \qquad
    \text{при }x\ne0.
\]
Локальную положительную определённость можно выразить оценкой
\[
    V(x)\ge\alpha_1(\lVert x\rVert),
    \qquad
    \alpha_1\in\mathcal K.
\]
Если она справедлива глобально с
\(\alpha_1\in\mathcal K_\infty\), то \(V\) радиально неограничена:
\[
    V(x)\to\infty
    \qquad
    \text{при }\lVert x\rVert\to\infty.
\]

Например,
\[
    V(x)=\frac12\lVert x\rVert^2
\]
положительно определена и радиально неограничена, тогда как
\[
    V(x)=1-e^{-\lVert x\rVert^2}
\]
положительно определена, но ограничена. Поэтому положительная
определённость и радиальная неограниченность являются различными
свойствами.

Функция \(W\) называется отрицательно определённой, если
\[
    W(0)=0,
    \qquad
    W(x)<0
    \quad
    \text{при }x\ne0.
\]
Это эквивалентно положительной определённости \(-W\). Отрицательная
полуопределённость означает \(W(x)\le0\). В прямом методе положительно
определённой выбирают саму функцию Ляпунова, а отрицательной или
полуотрицательной --- её производную вдоль траекторий.

\section{Прямой метод и экспоненциальный критерий}

Рассмотрим автономную систему
\[
    \dot x=g(x),
    \qquad
    g(0)=0,
\]
и непрерывно дифференцируемую функцию \(V\). Вдоль решения
\[
\begin{aligned}
    \dot V(x)
    &=
    \frac{\mathrm d}{\mathrm dt}V(x(t))
    \\
    &=
    \nabla V(x)^{\mathsf T}\dot x
    \\
    &=
    \nabla V(x)^{\mathsf T}g(x).
\end{aligned}
\]

\begin{definition}[Функция Ляпунова]
Положительно определённая непрерывно дифференцируемая функция \(V\)
называется функцией Ляпунова для равновесия \(x=0\), если
\[
    \dot V(x)\le0
\]
в некоторой его окрестности.
\end{definition}

Функция Ляпунова не обязана совпадать с физической энергией. Это
может быть специально сконструированная величина, удобная для
доказательства требуемого типа устойчивости.

\begin{theorem}[Прямой метод Ляпунова]
Пусть в окрестности начала координат существует непрерывно
дифференцируемая положительно определённая функция \(V\). Если
\[
    \dot V(x)\le0,
\]
то равновесие устойчиво по Ляпунову. Если
\[
    \dot V(x)<0
\]
для всех \(x\ne0\), то равновесие локально асимптотически устойчиво.
Если условия выполняются глобально и \(V\) радиально неограничена, то
вывод является глобальным.
\end{theorem}

Условие \(\dot V\le0\) не гарантирует притяжения: обобщённая энергия
может сохраняться на ненулевых траекториях. Строгое убывание вне
равновесия устраняет эту возможность.

\begin{theorem}[Квадратичный критерий]
Пусть существуют \(c_1,c_2,c_3>0\), для которых
\[
    c_1\lVert x\rVert^2
    \le
    V(x)
    \le
    c_2\lVert x\rVert^2
\]
и
\[
    \dot V(x)
    \le
    -c_3\lVert x\rVert^2.
\]
Тогда
\[
    \lVert x(t)\rVert
    \le
    \sqrt{\frac{c_2}{c_1}}
    \exp\left(
        -\frac{c_3}{2c_2}(t-t_0)
    \right)
    \lVert x(t_0)\rVert,
\]
то есть равновесие экспоненциально устойчиво.
\end{theorem}

\begin{proof}
Из верхней оценки на \(V\) следует
\[
    \lVert x\rVert^2\ge\frac1{c_2}V(x),
\]
поэтому
\[
    \dot V\le-\frac{c_3}{c_2}V.
\]
Интегрирование даёт
\[
    V(x(t))
    \le
    e^{-\frac{c_3}{c_2}(t-t_0)}V(x(t_0)).
\]
Используя обе квадратичные оценки, получаем
\[
\begin{aligned}
    c_1\lVert x(t)\rVert^2
    &\le
    V(x(t))
    \\
    &\le
    e^{-\frac{c_3}{c_2}(t-t_0)}
    c_2\lVert x(t_0)\rVert^2.
\end{aligned}
\]
После извлечения квадратного корня получается требуемое неравенство.
\end{proof}

Таким образом, в стандартной оценке можно взять
\[
    M=\sqrt{\frac{c_2}{c_1}},
    \qquad
    \alpha=\frac{c_3}{2c_2}.
\]

\section{Механическая энергия и принцип ЛаСалля}

Рассмотрим систему «масса--пружина--демпфер»
\[
    m\ddot q+b\dot q+kq=0,
    \qquad
    m,b,k>0.
\]
Положим
\[
    x=
    \begin{pmatrix}
        q\\
        v
    \end{pmatrix},
    \qquad
    v=\dot q.
\]
Тогда
\[
    \dot x
    =
    \begin{pmatrix}
        v\\[1mm]
        -\dfrac{k}{m}q-\dfrac{b}{m}v
    \end{pmatrix},
\]
а единственным равновесием является \((q,v)=(0,0)\).

Естественный кандидат --- полная механическая энергия
\[
    V_0(q,v)
    =
    \frac12kq^2+\frac12mv^2.
\]
Она положительно определена и радиально неограничена. Её производная
равна
\[
\begin{aligned}
    \dot V_0
    &=
    kqv+mv\dot v
    \\
    &=
    kqv+mv
    \left(
        -\frac{k}{m}q-\frac{b}{m}v
    \right)
    \\
    &=
    -bv^2
    \le0.
\end{aligned}
\]
Отсюда следует устойчивость по Ляпунову, но производная лишь
отрицательно полуопределена: при \(v=0\) и \(q\ne0\) она равна нулю.

Рассмотрим множество
\[
    \mathcal Z
    =
    \{(q,v):\dot V_0(q,v)=0\}
    =
    \{(q,v):v=0\}.
\]
Чтобы траектория оставалась в \(\mathcal Z\), необходимо
\(v(t)\equiv0\) и, следовательно, \(\dot v(t)\equiv0\). Но при \(v=0\)
\[
    \dot v=-\frac{k}{m}q,
\]
поэтому инвариантность требует \(q=0\). Наибольшее инвариантное
подмножество \(\mathcal Z\) состоит только из равновесия.

\begin{theorem}[Следствие принципа ЛаСалля]
Равновесие системы «масса--пружина--демпфер» глобально
асимптотически устойчиво.
\end{theorem}

Производная функции Ляпунова может обращаться в нуль в отдельных
неравновесных точках. Это не мешает сходимости, если траектория не
может навсегда остаться в соответствующем множестве.

\section{Перекрёстный член и строгая оценка}

Физическая энергия доказывает асимптотическую устойчивость, но не
даёт непосредственно строгой квадратичной оценки производной. Для
этого добавим малый перекрёстный член:
\[
    V_\varepsilon(q,v)
    =
    \frac12kq^2
    +
    \frac12mv^2
    +
    \varepsilon m qv,
\]
где \(\varepsilon>0\). В матричной форме
\[
    V_\varepsilon(x)
    =
    \frac12x^{\mathsf T}P_\varepsilon x,
    \qquad
    P_\varepsilon
    =
    \begin{pmatrix}
        k&\varepsilon m\\
        \varepsilon m&m
    \end{pmatrix}.
\]
Матрица \(P_\varepsilon\) положительно определена, если
\[
    0<\varepsilon<\sqrt{\frac{k}{m}}.
\]

Вдоль траекторий
\[
\begin{aligned}
    \dot V_\varepsilon
    ={}&
    kqv+mv\dot v
    +
    \varepsilon m(v^2+q\dot v)
    \\
    ={}&
    -bv^2+
    \varepsilon mv^2-
    \varepsilon bqv-
    \varepsilon kq^2.
\end{aligned}
\]
Следовательно,
\[
    \dot V_\varepsilon
    =
    -x^{\mathsf T}R_\varepsilon x,
\]
где
\[
    R_\varepsilon
    =
    \begin{pmatrix}
        \varepsilon k&\dfrac{\varepsilon b}{2}\\[2mm]
        \dfrac{\varepsilon b}{2}&b-\varepsilon m
    \end{pmatrix}.
\]
Для достаточно малого \(\varepsilon\) матрица \(R_\varepsilon\)
положительно определена. Достаточно обеспечить
\[
    b-\varepsilon m>0
\]
и
\[
    \varepsilon k(b-\varepsilon m)
    -
    \frac{\varepsilon^2b^2}{4}
    >0.
\]
Тогда существуют \(c_1,c_2,c_3>0\), такие что
\[
    c_1\lVert x\rVert^2
    \le
    V_\varepsilon(x)
    \le
    c_2\lVert x\rVert^2,
    \qquad
    \dot V_\varepsilon(x)
    \le
    -c_3\lVert x\rVert^2.
\]
По квадратичному критерию равновесие глобально экспоненциально
устойчиво.

Одна и та же система допускает разные функции Ляпунова. Энергия
\(V_0\) вместе с принципом ЛаСалля даёт асимптотическую устойчивость,
а модифицированная энергия \(V_\varepsilon\) --- непосредственную
экспоненциальную оценку. Поэтому слабый результат иногда отражает не
свойство системы, а недостаточно точный выбор функции Ляпунова.

\section{Градиентный поток и три естественных кандидата}

Рассмотрим задачу
\[
    \min_{x\in\RR^n}f(x)
\]
и градиентный поток
\[
    \dot x=-\nabla f(x).
\]
Стационарная точка \(x^\star\), удовлетворяющая
\[
    \nabla f(x^\star)=0,
\]
является равновесием. Для выпуклой функции она является глобальным
минимумом, а для строго или сильно выпуклой функции минимум
единственен.

Пусть \(f\in C^2\) и
\[
    \nabla^2f(x)\succeq\mu I,
    \qquad
    \mu>0.
\]
Для функции
\[
    V_\nabla(x)
    =
    \frac12\lVert\nabla f(x)\rVert^2
\]
имеем
\[
\begin{aligned}
    \dot V_\nabla
    &=
    \nabla f(x)^{\mathsf T}\nabla^2f(x)\dot x
    \\
    &=
    -\nabla f(x)^{\mathsf T}
    \nabla^2f(x)\nabla f(x)
    \\
    &\le
    -\mu\lVert\nabla f(x)\rVert^2
    \\
    &=
    -2\mu V_\nabla.
\end{aligned}
\]
Следовательно,
\[
    \lVert\nabla f(x(t))\rVert
    \le
    e^{-\mu(t-t_0)}
    \lVert\nabla f(x(t_0))\rVert.
\]

Сильная выпуклость также даёт
\[
    \langle
        \nabla f(x)-\nabla f(x^\star),
        x-x^\star
    \rangle
    \ge
    \mu\lVert x-x^\star\rVert^2.
\]
Поскольку \(\nabla f(x^\star)=0\), получаем
\[
    \lVert\nabla f(x)\rVert
    \ge
    \mu\lVert x-x^\star\rVert.
\]
Если градиент \(L\)-липшицев, то
\[
    \lVert\nabla f(x(t_0))\rVert
    \le
    L\lVert x(t_0)-x^\star\rVert,
\]
и потому
\[
    \lVert x(t)-x^\star\rVert
    \le
    \frac{L}{\mu}
    e^{-\mu(t-t_0)}
    \lVert x(t_0)-x^\star\rVert.
\]

Более прямой выбор ---
\[
    V_x(x)
    =
    \frac12\lVert x-x^\star\rVert^2.
\]
Тогда
\[
\begin{aligned}
    \dot V_x
    &=
    -\langle x-x^\star,\nabla f(x)\rangle
    \\
    &\le
    -\mu\lVert x-x^\star\rVert^2
    \\
    &=
    -2\mu V_x.
\end{aligned}
\]
Следовательно,
\[
    \lVert x(t)-x^\star\rVert
    \le
    e^{-\mu(t-t_0)}
    \lVert x(t_0)-x^\star\rVert.
\]

Наконец, для ошибки по значению
\[
    V_f(x)=f(x)-f^\star
\]
всегда выполняется
\[
    \dot V_f
    =
    -\lVert\nabla f(x)\rVert^2.
\]
Если функция удовлетворяет неравенству Поляка--Лоясиевича
\[
    \frac12\lVert\nabla f(x)\rVert^2
    \ge
    \mu\bigl(f(x)-f^\star\bigr),
\]
то
\[
    \dot V_f\le-2\mu V_f
\]
и
\[
    f(x(t))-f^\star
    \le
    e^{-2\mu(t-t_0)}
    \bigl(f(x(t_0))-f^\star\bigr).
\]
PL-условие тем самым обеспечивает экспоненциальное убывание ошибки
по значению даже без выпуклости.

\section{Множество решений и выбор функции}

PL-неравенство не гарантирует единственности минимизатора. Например,
\[
    f(x_1,x_2)=\frac12x_1^2
\]
удовлетворяет
\[
    \frac12\lVert\nabla f(x)\rVert^2
    =
    f(x)-f^\star,
\]
но
\[
    \mathcal X^\star
    =
    \{(0,x_2):x_2\in\RR\}.
\]
В этом случае \(V_f=f-f^\star\) положительно определена относительно
множества решений:
\[
    V_f(x)=0
    \quad\Longleftrightarrow\quad
    x\in\mathcal X^\star.
\]
Экспоненциальная оценка гарантирует сходимость значения функции, но
сама по себе не выбирает единственный элемент множества
\(\mathcal X^\star\). Верное общее следствие PL-свойства состоит в том,
что каждая стационарная точка является глобальным минимумом.

Три основных кандидата отвечают разным целям. Функция
\[
    \frac12\lVert x-x^\star\rVert^2
\]
естественна при сильной выпуклости и непосредственно оценивает
расстояние до решения. Функция
\[
    f(x)-f^\star
\]
сочетается с PL-неравенством и контролирует ошибку по значению.
Функция
\[
    \frac12\lVert\nabla f(x)\rVert^2
\]
удобна при положительной определённости гессиана и даёт оценку нормы
градиента.

Для инерционных алгоритмов функция Ляпунова обычно включает ошибку по
функции, квадратичный член по скорости и перекрёстные слагаемые. Для
распределённых алгоритмов к ним добавляется ошибка согласования,
например
\[
    \sum_{i=1}^{N}\lVert x_i-\bar x\rVert^2.
\]
Прямой метод не требует явного решения системы: достаточно найти
равновесие, построить положительно определённую функцию, вычислить её
производную и оценить скорость убывания. Основная трудность состоит
в выборе \(V\). Обратные теоремы Ляпунова часто гарантируют
существование подходящей функции, но не дают автоматически её простой
аналитической формы.
